\subsection{Оптические гироскопы}

Оптические гироскопы не имеют подвижных механических частей в традиционном понимании. Их работа основана на явлении интерференции электромагнитных волн и эффекте Саньяка \cite{malkin_2002}.

Эффект Саньяка заключается в следующем: два когерентных луча света, распространяющиеся по одному и тому же замкнутому оптическому контуру в противоположных направлениях, приобретают разность фаз $\Delta \Phi$ при вращении контура с угловой скоростью $\Omega$:

\begin{equation}
	\Delta \Phi = \frac{8 \pi S}{\lambda c} \Omega
\end{equation}

где S — площадь, охватываемая контуром, $\lambda$ - длина волны света, c — скорость света.

% -----------------------------------------------------------------------------

\subsubsection{Лазерный гироскоп}

Конструктивные особенности лазерного гироскопа определяются необходимостью создания сверхстабильного кольцевого лазерного резонатора: его основу составляет цельный монолитный блок из сверхстабильной стеклокерамики (например, Церодур), в котором фрезерованы каналы, образующие замкнутый треугольный или квадратный контур. В углах резонатора установлены высокоточные многослойные диэлектрические зеркала, формирующие кольцевой оптический контур, причём одно из зеркал является полупрозрачным для вывода излучения на фотодетектор \cite{opticheskie_giroskopy_2005}. 

При вращении гироскопа резонансные частоты для встречных волн становятся разными. Возникает разность частот $\Delta f$ (частотный сдвиг Саньяка \cite{malkin_2002}), которая строго пропорциональна угловой скорости вращения $\Omega$. 

\begin{equation}
	\Delta f = \frac{4 S}{\lambda L} \Omega
\end{equation}

где L - периметр контура.

Таким образом, измеряя разность частот $\Delta f$, можно напрямую и с высокой точностью определить угловую скорость $\Omega$.

% -----------------------------------------------------------------------------

\subsubsection{Волоконно-оптический гироскоп}

Волоконно-оптический гироскоп — это тип оптического гироскопа, в котором для усиления эффекта Саньяка используется многовитковая катушка из оптического волокна \cite{volokonno_opticheskiy_1987}. Это пассивный интерферометрический прибор, ставший основной альтернативой и конкурентом лазерного гироскопа благодаря своей конструктивной простоте и высокой надежности.

Свет от внешнего широкополосного источника разделяется надвое, и два луча запускаются в противоположных направлениях в одну и ту же волоконную катушку. После прохождения контура лучи снова сводятся и интерферируют.

При вращении гироскопа возникает разность фаз $\Delta \Phi$ между встречными лучами, которая регистрируется фотодетектором \cite{volokonno_opticheskiy_1987}:

\begin{equation}
	\Delta \Phi = \frac{2 \pi L D}{\lambda c} \Omega
\end{equation}

где L — длина волокна, D — диаметр катушки, $\lambda$ — длина волны света, c — скорость света, $\Omega$ — угловая скорость.

Ключевой конструктивной особенностью волоконно-оптического гироскопа является его модульная структура, центральным элементом которой служит многовитковая катушка (длиной от сотен метров до нескольких километров) из одномодового, сохраняющего поляризацию оптического волокна, намотанная особым способом для минимизации температурных и вибрационных помех. Весь оптический тракт формируется вокруг интегрально-оптического чипа, выполняющего функции разделителя, поляризатора и фазового модулятора, а в качестве источника излучения используется широкополосный суперлюминесцентный диод для подавления паразитных интерференционных шумов.

